% TOP_PROBLEM: {"id": 30006897, "title": "Lieb-Thirring Conjecture on the Sharp Constant for Riesz Means of Order One in Three Dimensions", "category_id": 8, "category": {"id": 8, "name": "analysis", "display_name": "Analysis", "description": "Limits, continuity, calculus, and function theory.", "slug": "analysis", "order_index": 8, "created_at": "2026-07-31T15:26:25.671Z"}, "status": "open"}
% FIRST_POSTED: 2026-09-27
% !TeX program = lualatex
% ATTEMPT_STATUS: unresolved
% Model: GPT-6 Astra Ultra; continuation dated 27 September 2026.
\documentclass[11pt]{article}
\usepackage[margin=25mm]{geometry}
\usepackage{fontspec}
\setmainfont{Latin Modern Roman}
\usepackage{amsmath,amssymb,mathtools}
\usepackage{unicode-math}
\setmathfont{Latin Modern Math}
\usepackage{xurl,hyperref}
\hypersetup{colorlinks=true,urlcolor=blue}
\setcounter{secnumdepth}{0}
\setlength{\parindent}{0pt}
\setlength{\parskip}{5pt}
\setlength{\emergencystretch}{3em}
\begin{document}
\section*{\texorpdfstring{Lieb-Thirring Conjecture on the Sharp Constant for Riesz Means of Order One in Three Dimensions}{Lieb-Thirring Conjecture on the Sharp Constant for Riesz Means of Order One in Three Dimensions}}\label{30006897}
\subsection{Definitions and mathematical statement}
For a real potential $V$ on ${\mathbb{R}}^3$, let $V_- =\max\{-V,0\}$ and form the self-adjoint Schr\"odinger operator $H=-\Delta+V$ through its quadratic form. Its negative spectral part is $H_-=\max\{-H,0\}$, so $\operatorname{Tr}H_-$ is the sum of magnitudes of negative eigenvalues, with multiplicity. The assertion is
\[
 \operatorname{Tr}(-\Delta+V)_-\le
 \frac1{15\pi^2}\int_{{\mathbb{R}}^3}V_-^{5/2}{\,\mathrm{d}} x,
\]
for the admissible potential class of the Lieb--Thirring inequality, with this constant optimal. It suffices for the constant question to formulate the supremum over smooth compactly supported negative potentials and use the standard form approximation extension. The exponent is the order-one, three-dimensional case; other Riesz orders or dimensions have different sharp constants.
\par\smallskip\textit{Formulation and concept references:} \href{https://arxiv.org/abs/2007.09326}{[S1]}.

\subsection{Short English statement}
Is the semiclassical value $1/(15\pi^2)$ the exact best constant controlling total negative Schr\"odinger energy by the integral of the potential's negative part to power $5/2$?
\subsection{Research attempt}
\paragraph{Dimension and current source scope.}
The target is $\gamma=1$ in dimension three, with one scalar
Schrodinger species and kinetic operator $-\Delta$. The factors
below correspond to this convention. The survey [S1] is an account
of the sharp-constant problem, not a proof of this case. A newer
primary preprint [E1], version 2 of 14 September 2026, claims the
sharp one-dimensional order-one inequality and an operator-valued
extension giving $L_{1,d}\leq(2/\sqrt3)L_{1,d}^{\rm cl}$. That
stated higher-dimensional factor is strictly larger than one.
Even accepting the manuscript's results does not identify the
three-dimensional constant requested here. This audit verifies its
stated scope; it does not independently validate its full proof.

We establish necessity of the printed constant by explicit box
trials, derive the equivalent kinetic-energy inequality with all
constants, and examine two finite-rank tests. The unproved step is
the matching kinetic lower bound for arbitrary orthonormal families.

\paragraph{Reduction to negative potentials and scaling.}
Writing $V=V_+-V_-$, quadratic-form ordering gives
$-\Delta+V\geq-\Delta-V_-$. The variational characterization of
negative energy therefore reduces the upper-bound question to
operators $-\Delta-W$ with $W\geq0$. For $\lambda>0$, the unitary
dilation $(U_\lambda\psi)(x)=\lambda^{3/2}\psi(\lambda x)$ gives
\[
 -\Delta-\lambda^2W(\lambda x)
     =\lambda^2U_\lambda(-\Delta-W)U_\lambda^{-1}.
\]
Both the negative trace and $\int(\lambda^2W(\lambda x))^{5/2}dx$
scale by $\lambda^2$. Thus rescaling a single potential cannot
improve its ratio. Spatial concentration is not a way around a
nonsharp estimate that is already scale invariant.

For any finite orthonormal family $\psi_j\in H^1(\mathbb R^3)$,
with occupation numbers $0\leq a_j\leq1$, the variational principle
gives
\[
 \operatorname{Tr}(-\Delta-W)_-
 \geq\sum_j a_j\left(\int W|\psi_j|^2-\int|\nabla\psi_j|^2\right).
 \tag{386.1}
\]
Taking the supremum over such finite-rank positive contractions
recovers the negative trace. This follows first from the spectral
theorem by occupying the negative spectral subspace, and then by
monotone finite-rank approximation when the trace is not a finite
sum. Trial subspaces produce lower bounds; they do not prove the
desired upper bound.

\paragraph{The semiclassical constant is forced by box trials.}
Let $Q_R=(0,R)^3$ and choose a smooth compactly supported $W_R$
with $0\leq W_R\leq1$, equal to one on $Q_R$, and supported in
$(-1,R+1)^3$. Thus
\[
 \int W_R^{5/2}=R^3+O(R^2).
\]
Use the orthonormal Dirichlet eigenfunctions of the cube, extended
by zero outside it. They belong to $H^1(\mathbb R^3)$, and their
kinetic eigenvalues are $\pi^2|m|^2/R^2$ for
$m\in\mathbb N^3$, where the positive integers start at one.
Selecting those with eigenvalue below one in (386.1) gives
\[
 \operatorname{Tr}(-\Delta-W_R)_-
 \geq\sum_{m\in\mathbb N^3}
       \left(1-\frac{\pi^2|m|^2}{R^2}\right)_+.
 \tag{386.2}
\]
Divide by $R^3$. This is a Riemann sum for a compactly supported
continuous function on the positive octant, so its limit is
\[
 \frac1{\pi^3}\int_{\substack{\xi_i\geq0\\|\xi|<1}}
                  (1-|\xi|^2)\,d\xi
 =\frac1{8\pi^3}\,4\pi\int_0^1(r^2-r^4)\,dr
 =\frac1{15\pi^2}.
 \tag{386.3}
\]
The boundary layer in the potential has negligible relative volume.
Consequently any constant valid for the smooth compactly supported
potential class must be at least $1/(15\pi^2)$. No appeal to a
formal phase-space approximation is required for this lower bound.
The Riemann-sum limit and the trial-space inequality provide a
rigorous sequence of admissible tests.

This does not say that every negative eigenfunction resembles a
cube sine wave, or that an arbitrary potential has no additional
quantum contribution. Those would be upper-bound statements, which
the variational trials cannot establish.

\paragraph{The exactly dual kinetic formulation.}
For a finite-rank contraction
$\gamma=\sum_j a_j|\psi_j\rangle\langle\psi_j|$, put
\[
 \rho_\gamma(x)=\sum_j a_j|\psi_j(x)|^2,\qquad
 T(\gamma)=\sum_j a_j\int|\nabla\psi_j|^2.
\]
The sharp potential inequality with constant $L$ implies
\[
 T(\gamma)\geq K(L)\int\rho_\gamma^{5/3},
 \qquad K(L)=\frac35\left(\frac2{5L}\right)^{2/3}.
 \tag{386.4}
\]
To prove this, (386.1) and the potential upper bound give
$T(\gamma)\geq\int(W\rho_\gamma-LW^{5/2})$ for every nonnegative
test potential. For each fixed density value $\rho$, the maximizing
$W$ solves $\rho=(5L/2)W^{3/2}$ and has value
$K(L)\rho^{5/3}$. Truncation and smooth approximation of the
maximizing potential justify the pointwise optimization under the
integral, initially for smooth finite families and then by the
usual form approximation. A bounded compactly supported truncated
potential can be approximated in the required integral norms, so
no irregular potential is silently added to the initial test class.

Conversely, if (386.4) holds for all such $\gamma$ with a constant
$K$, then
\[
 \int W\rho_\gamma-T(\gamma)
 \leq\int\sup_{\rho\geq0}(W\rho-K\rho^{5/3})
 =\frac25\left(\frac3{5K}\right)^{3/2}\int W^{5/2}.
\]
Taking the variational supremum in (386.1) proves the potential
inequality with
\[
 L(K)=\frac25\left(\frac3{5K}\right)^{3/2}.
 \tag{386.5}
\]
The formulas (386.4)--(386.5) are mutual inverses. At the printed
semiclassical constant they give
\[
 K_{\rm cl}=\frac35(6\pi^2)^{2/3}.
 \tag{386.6}
\]
Thus the exact remaining inequality is
\[
 \boxed{\ T(\gamma)\geq\frac35(6\pi^2)^{2/3}
                    \int\rho_\gamma^{5/3},\qquad0\leq\gamma\leq1.\ }
 \tag{386.7}
\]
The constraint $\gamma\leq1$ is the orthogonality/occupation
restriction. Dropping it produces a different, false universal
inequality with this density scaling.

\paragraph{Why a one-function inequality is not enough.}
Suppose one proves $\int|\nabla\psi|^2\geq C\int|\psi|^{10/3}$
for normalized single functions. Summing over an orthonormal family
controls $\sum_j|\psi_j|^{10/3}$, whereas (386.7) contains
$(\sum_j|\psi_j|^2)^{5/3}$. For positive numbers,
\[
 \left(\sum_{j=1}^N b_j\right)^{5/3}
       \leq N^{2/3}\sum_{j=1}^N b_j^{5/3},
\]
and a factor of this size is unavoidable when the $b_j$ are equal.
Orthogonality is global in space, not pointwise disjointness of the
densities. A proof that sums scalar inequalities therefore loses
precisely the particle-number independence required here unless it
uses additional information.

If occupation numbers are unconstrained, repeating one normalized
state $N$ times makes kinetic energy proportional to $N$ and the
density integral proportional to $N^{5/3}$, directly disproving a
uniform kinetic inequality. This is a countercheck on the role of
the fermionic constraint, not an admissible counterexample to the
conjecture.

\paragraph{A completely evaluated rank-one test.}
Let
\[
 \psi_\sigma(x)=(\pi\sigma^2)^{-3/4}
                 \exp(-|x|^2/(2\sigma^2)).
\]
Gaussian integration gives
\[
 \|\psi_\sigma\|_2=1,\qquad
 T=\frac3{2\sigma^2},\qquad
 \int|\psi_\sigma|^{10/3}
       =\frac1{\pi\sigma^2}\left(\frac35\right)^{3/2}.
\]
Its kinetic ratio is therefore
\[
 \frac{T}{\int\rho^{5/3}}
       =\frac{3\pi}{2}\left(\frac53\right)^{3/2},
 \tag{386.8}
\]
independent of $\sigma$. This value is larger than $K_{\rm cl}$;
squaring their positive ratio reduces the comparison to an elementary
numerical inequality in $\pi$. Thus this Gaussian does not violate
(386.7), and it does not saturate it. Testing an entire dilation
family does not add independent evidence because the ratio is
scale invariant.

The same caution applies to widely separated wells. If compactly
supported finite families are translated so that their supports
are disjoint, both their kinetic energies and the integrals of
their total densities to power $5/3$ add exactly. Replicating the
same family in disjoint regions preserves its ratio. An improved
ratio would require a genuinely new overlapping or many-state
structure, not repetition of an existing test.

\paragraph{The local-density heuristic and its unproved step.}
At a fixed point, filling momentum states inside a ball produces
\[
 \rho=(2\pi)^{-3}\frac{4\pi}{3}k_F^3,\qquad
 t=(2\pi)^{-3}\frac{4\pi}{5}k_F^5
       =K_{\rm cl}\rho^{5/3}.
\]
This explains the constant and agrees with the box construction.
It is not a pointwise inequality for a general density matrix:
position and momentum localization cannot both be exact, and a
phase-space representation of a quantum state need not be a
nonnegative occupation function bounded by one. Replacing it by
such a function without an error estimate is the first invalid
step in a common heuristic proof.

One can localize into boxes and compare with Dirichlet or Neumann
energies, but localization derivatives and boundary terms appear.
For the lower sharpness sequence their relative effect vanishes as
the box grows. For an arbitrary inhomogeneous potential there is
no single box scale at which those errors vanish uniformly while
retaining the exact proposed constant. Keeping a positive error
term yields a nonsharp bound; dropping it is not justified.

\paragraph{Diagnostics and outcome.}
The finite script checks the Legendre constants, evaluates the
Gaussian ratio and the semiclassical value, and computes Dirichlet
box sums along an increasing sequence of $R$. The convergence of
the sums is proved in (386.3), not inferred from finitely many data
points. These tests do not search all potentials or all finite-rank
contractions.

The proposed constant is rigorously necessary, and (386.7) is an
equivalent exact target for sufficiency. The recent one-dimensional
claim has been kept separate from this three-dimensional question.
The first unresolved step is proving (386.7) with no factor loss
for every admissible orthonormal family. This attempt does not
supply that step and does not settle the printed conjecture.
\subsection{Sources}
\begin{itemize}
\item[S1]Rupert L. Frank, The Lieb-Thirring inequalities: Recent results and open problems, arXiv:2007.09326 (2020).
\url{https://arxiv.org/abs/2007.09326}.
\textit{Formulation source.}
\item[E1]Larry Read and Marvin R. Schulz, \emph{The sharp one-dimensional
Lieb--Thirring inequality for the sum of eigenvalues},
arXiv:2609.10478v2, 14 September 2026.
\url{https://arxiv.org/abs/2609.10478v2}.
\textit{Primary abstract inspected 27 September 2026. It claims the sharp
one-dimensional result and the higher-dimensional bound with factor
$2/\sqrt3$; neither assertion is substituted for a sharp result in
dimension three. The full manuscript proof was not independently checked.}
\end{itemize}
\end{document}
